\section{Google DeepMind、世界モデルGenie 4を発表}
Google DeepMindは、操作可能な世界を生成する「Genie 4」を研究プレビューとして発表した。自然言語によるリアルタイム操作、環境音・音楽・対話のネイティブ生成、Geminiとの統合による整合性向上を特徴として掲げる。公式説明では、物理表現の改善と最大約2分の一貫したシーン生成も主張する。まず信頼されたテスターに提供し、Project GenieやAPI等への展開を予定するという。ロボット学習などへの応用が期待されるが、性能や利用条件の独立評価はこれからである。
\dailyxsource{Google DeepMind公式}{https://x.com/GoogleDeepMind/status/2115114028338962672}

\section{Liquid AI、小規模マルチモーダルLFM3を公開}
Liquid AIは、約17億パラメータの共通基盤から派生したオープンウェイトの「LFM3」ファミリーを発表した。テキスト、視覚、音声、複合入力、意思決定などを扱う計6モデルを示し、長文コンテキストや単一GPUでのファインチューニングを訴求する。公式ブログでは畳み込み、スライディングウィンドウ、グローバル注意機構を組み合わせた構成と、22兆トークンの事前学習を説明している。性能値は提供元の評価に基づき、ライセンスや実機での速度も確認が必要だ。
\dailyxsource{コミュニティ投稿（AJ）}{https://x.com/ItsmeAjayKV/status/2115146675648131448}

\section{OpenAIとAnthropic、ETH反証を主張}
OpenAIとAnthropicが、それぞれ内部AIモデルによって計算量理論の指数時間仮説（ETH）を反証したとする発表・論文がXで議論された。Leanによる形式化の公開や専門家グループの検証計画が紹介される一方、両社の主張が数学的に成立するかは\textbf{独立検証が未完了}である。専門家の関与も証明の承認を意味しない。OpenAI側の公式発表は今回の観測窓外にあり、本号が扱うのは主に窓内の反応である。結果や計算量への影響を確定事実として受け取る段階ではない。

\section{Tencent Hunyuan、Hy-Nomasを公開}
Tencent Hunyuanの「Hy-Nomas」が、テキスト指示とカメラ操作によって動画世界を生成するモデルとして公開された。ModelScopeの投稿ではApache 2.0ライセンス、8ステップ蒸留、量子化やCPUオフロードを利用した単一GPU推論への対応を紹介する。コミュニティでは低VRAM環境での実行報告も見られるが、リアルタイム性能が確認されたわけではない。Genie 4と同じ世界モデル領域で、公開・ローカル実行可能性を重視する選択肢として話題になった。
\dailyxsource{ModelScope公式}{https://x.com/ModelScope2022/status/2114905680171745321}

\section{Cloudflare、推論時の顧客データ露出を開示}
Cloudflareは、Workers AIで提供していたGLM-5.3の長文処理に関連し、別顧客のプロンプト断片が応答へ混入した過去のインシデントを開示した。Grokの収集によれば、影響期間は8月27日から9月9日で、155件が対象と報告されている。同社は修正済みであり、意図的な悪用を示す証拠はないと説明する。KVキャッシュ周辺のテナント間隔離が、モデル品質だけでなく情報保護にも直結することを示す事例として注目された。
\dailyxsource{Gadi Evronによる関連投稿}{https://x.com/gadievron/status/2114895196664238484}

\section{AnthropicのS-1提出、財務面の分析が続く}
Anthropicの新規株式公開に向けたS-1登録届出書について、観測窓内では収益成長、計算資源への支出コミットメント、効率指標などの分析投稿が続いた。今回の話題は新たな公式提出の発表ではなく、既報を受けた市場・コミュニティの解釈が中心である。公式発表自体は観測窓外であり、個別の財務数値も投稿者の説明や報道に依存する。事業規模や資金需要を判断するには、正式なSEC文書の注記を確認する必要がある。

\section{オープンモデルの価格性能とローカル実行を比較}
DeepSeek V4.1 Flash、GLM系、LFM3などのオープンウェイトモデルを比較する投稿が複数共有された。速度、費用、品質、手元のGPUで実行できるかという観点が交錯し、主観的な順位表や実機報告が話題となった。公式ベンチマークとコミュニティの体験談が混在しているため、単一の順位を客観的な結論とみなすことはできない。利用環境と評価条件をそろえた比較が、今後のモデル選択には重要となる。
\dailyxsource{Miaによる比較投稿}{https://x.com/MiaAI_lab/status/2115201177784901788}

\section{世界モデル・エージェントの周辺動向}
Genie 4などの世界モデル発表を受け、関連する買収報道やエージェント応用への期待がX上で散発的に言及された。今回のGrok収集では、これらは二次的な投稿が中心で、独立した主要ニュースとして扱うだけの一次情報は確認されていない。世界の状態を継続的に扱うモデルと、外部環境で行動するエージェントの接点が関心を集めているものの、具体的な成果については続報を待つ必要がある。
